\documentclass[11pt]{article}
\usepackage{putnamprep}

\begin{document}
\sessiontitle{Session 3: Solutions}{Do not open until your write-up is done.}

\section*{Check questions from the lesson}

\paragraph{Finite subsets.} $\{x_3, x_7, x_8\}$ has largest index $8$, so it is in batch 8. Batch 4 has $2^3 = 8$ subsets:
\[ \{x_4\},\ \{x_1, x_4\},\ \{x_2, x_4\},\ \{x_3, x_4\},\ \{x_1, x_2, x_4\},\ \{x_1, x_3, x_4\},\ \{x_2, x_3, x_4\},\ \{x_1, x_2, x_3, x_4\}. \]

\paragraph{Identities.} For (a) with $n = 3$:
\[ (a - b)(a^2 + ab + b^2) = a^3 + a^2 b + ab^2 - a^2 b - ab^2 - b^3 = a^3 - b^3. \]
For (e), with $A = x^2 + 2y^2$ and $B = 2xy$:
\[ (A + B)(A - B) = A^2 - B^2 = x^4 + 4x^2y^2 + 4y^4 - 4x^2y^2 = x^4 + 4y^4. \]

\paragraph{Mod 4.} Yes. $21 = 3 \cdot 7$ is $\equiv 1 \pmod 4$, and both of its prime factors are $\equiv 3 \pmod 4$.

\paragraph{Euclid on the list $2, 3, 5$.} $N = 2 \cdot 3 \cdot 5 + 1 = 31$, which is prime. So the prime dividing $N$ is $31$, and it is not on the list.

\section*{Three quick checks}

\paragraph{Exercise 1.} With numbers: $63 = 4^3 - 1 = (4 - 1)(4^2 + 4 + 1) = 3 \cdot 21$.

In general, let $n$ be composite and write $n = rs$ with $1 < r, s < n$. Apply identity (a) to $2^r$ and $1$ with exponent $s$:
\[ 2^n - 1 = (2^r)^s - 1 = (2^r - 1)\left((2^r)^{s-1} + (2^r)^{s-2} + \cdots + 1\right). \]
Since $1 < r < n$, the first factor satisfies $1 < 2^r - 1 < 2^n - 1$. So $2^n - 1$ has a divisor strictly between $1$ and itself, and is composite.

\paragraph{Exercise 2.} With $y = 1$, identity (e) gives $n^4 + 4 = (n^2 + 2n + 2)(n^2 - 2n + 2)$. With numbers: $20 = 10 \cdot 2$ and $85 = 17 \cdot 5$.

The smaller factor is $n^2 - 2n + 2 = (n-1)^2 + 1$. For $n \ge 2$ it is at least $1^2 + 1 = 2$. The larger factor is bigger still. So both factors are at least $2$, and $n^4 + 4$ is not prime.

\paragraph{Exercise 3.} The step fails at $n = 1$, that is, going from one horse to two. Remove horse A: what remains is $\{B\}$. Remove horse B: what remains is $\{A\}$. Those two groups share no horse, so nothing links A's color to B's. The step silently needs a horse that stays in both groups, which requires $n \ge 2$. So the basis would have to include $n = 2$ as well, and ``any two horses are the same color'' is false. This is the two-base-case trap: the basis did not reach as far as the step needed.

\section*{The main proof}

\paragraph{Why ``product plus one'' fails.} For the pretend list $\{3\}$ it gives $4 = 2 \cdot 2$, which has no prime factor of the form $4k - 1$. For $\{3, 7, 11\}$ it gives $232 = 2^3 \cdot 29$, and $29 \equiv 1 \pmod 4$. Euclid's number gives \emph{a} new prime, but not one of the kind we need.

\begin{claim}
There are infinitely many primes of the form $4k - 1$, that is, primes congruent to $3 \pmod 4$.
\end{claim}
\begin{proof}
Suppose, for contradiction, that there are only finitely many such primes, and let $p_1, p_2, \dots, p_n$ be all of them. The list is not empty, since $3$ is such a prime. Let
\[ d = 4\,p_1 p_2 \cdots p_n - 1. \]
\begin{enumerate}
  \item $d \ge 4 \cdot 3 - 1 = 11 > 1$, so $d$ is a product of one or more primes.
  \item $d$ is odd, since $4 p_1 \cdots p_n$ is even. So $2$ is not a prime factor of $d$.
  \item No $p_i$ divides $d$. Indeed $p_i$ divides $4 p_1 \cdots p_n$, so if $p_i$ also divided $d$ it would divide their difference, which is $1$. No prime divides $1$.
  \item By steps 2 and 3, every prime factor of $d$ is odd and is not one of $p_1, \dots, p_n$. An odd prime is congruent to $1$ or $3 \pmod 4$, and $p_1, \dots, p_n$ are all the primes congruent to $3$. So every prime factor of $d$ is congruent to $1 \pmod 4$.
  \item A product of integers congruent to $1 \pmod 4$ is congruent to $1 \pmod 4$, because $(4s+1)(4t+1) = 4(4st + s + t) + 1$. So $d \equiv 1 \pmod 4$.
  \item But $d = 4(p_1 \cdots p_n) - 1 = 4(p_1 \cdots p_n - 1) + 3$, so $d \equiv 3 \pmod 4$.
\end{enumerate}
Steps 5 and 6 give $d$ two different remainders mod $4$. This contradiction shows that there are infinitely many primes of the form $4k - 1$.
\end{proof}

\begin{gradingnote}\leavevmode
\begin{itemize}
  \item Steps 2 and 5 must both appear. ``$d$ is odd'' rules out the prime $2$, which is neither $1$ nor $3 \pmod 4$. Skipping either one costs points.
  \item Step 3 needs a reason. ``$d$ is not divisible by any $p_i$'' with no justification is a gap.
  \item Step 1 is one line and easy to forget. Without it, ``every prime factor of $d$'' might be talking about nothing.
  \item You do not need $d$ to be prime, and it usually is not: for the list $\{3, 7, 11\}$, $d = 923 = 13 \cdot 71$. The proof never claims $d$ is prime. A write-up that says ``so $d$ is a new prime'' contains a false claim.
\end{itemize}
\end{gradingnote}

\paragraph{The check question you kept.} Why does the same proof not work for primes of the form $4k + 1$? Because the Lemma has no mirror image. A product of numbers $\equiv 1$ must be $\equiv 1$, which is what forced a factor $\equiv 3$ above. But a number $\equiv 1$ can be built entirely from primes $\equiv 3$, such as $21 = 3 \cdot 7$. So knowing a number is $\equiv 1 \pmod 4$ does not force it to have a prime factor $\equiv 1$. The $4k+1$ case is true, but it needs a different idea, which comes later.
\end{document}
